A. QUALITATIVE SYLLOGISM

The first syllogism is the syllogism of existence or the qualitative syllogism,
as it was portrayed in the previous section, (I) I - P - U
[individuality, particularity, universality]
that a subject as individual is joined together, through a quality,
with some universal determinacy.

    That the subject (terminus minor) has even further determinations than
    that of individuality, similarly that the other extreme
    (the predicate of the conclusion, the terminus maior) is
    further determined than being merely a universal,
    does not come into consideration here;
    only the forms through which they constitute the syllogism [come into consideration].

Addition. The syllogism of existence is a syllogistic inference merely at the level of
understanding and, indeed, insofar as the individuality, the particularity, and the
universality stand opposite one another in an entirely abstract manner here. Thus,
this syllogism is then the most extreme way that the concept comes to be outside
itself. We have here something immediately individual as a subject; some particular
side, a property, in this subject is then emphasized and by means of this property
the individual demonstrates itself to be a universal. So, for example, we say 'this
rose is red; red is a colour, therefore, this rose is something coloured'. It is this form
[Gestalt] of the syllogism, above all, that is typically discussed in ordinary logic. In
former times, the syllogism was considered the absolute rule of all knowing and a
scientific claim obtained then as something justified only if it was demonstrated
in a manner mediated by a syllogism. Today one encounters the diverse forms
of the syllogism almost exclusively only in compendia of logic, and acquaintance
with those various forms counts as empty pedantry, of no further use of any sort
either in practical life or even in science. In this regard, it deserves to be noted,
first, that although it would be superfluous and pedantic to enter on the scene at
each occasion with the entire elaboration of formal modes of inferring, the diverse
forms of inference nonetheless continue to impose themselves on our knowing.
For example, if someone waking up in the morning during the wintertime hears
carriages clanging on the street and this occasions him to consider that things
may well have frozen solid, he performs an operation of inferring and we repeat
this operation daily amidst the most manifold complications. Becoming explicitly
conscious of this, one's daily actions as a thinking being might thus at least be
of no slighter interest than the well-recognized interest in becoming acquainted,
not only with the functions of our organic life, e.g. the functions of digestion,
production of the blood, breathing, and so forth, but also with the functions of
the processes and formations of nature surrounding us. It will undoubtedly have
to be conceded in this connection that just as little as a foregoing study of anatomy
and physiology is needed in order to digest properly, to breath properly, and so
on, one needs to have studied logic first in order to draw proper conclusions. - It
is Aristotle who first observed and described the diverse forms and so-called figures
of the syllogism in their subjective meaning and, indeed, did so with such sureness
and determinacy that essentially nothing further had to be added. Although this
accomplishment brings Aristotle great honour, by no means is it the forms of
syllogistic inference at the level of understanding or at the level generally of finite
thinking that he employed in his genuine philosophical investigations.

This syllogism is (a) completely contingent with respect to its determi-
nations since the middle, as an abstract particularity, is merely any sort of
determinacy of the subject, of which, as something immediate and thus
empirically concrete, it has several. Hence, it can be joined together just as
much with many sorts of other universalities, just as an individual particu-
larity in turn can also have several diverse determinacies in itsel£ Thus, the
subject can be related to different universals by means of the same medius
terminus [middle term].

    It is more that formally inferring has gone out of fashion than that
    its incorrectness has been detected and that the lack of its use has
    been justified on that basis [i.e. its incorrectness). This and the
    following section indicate the vacuousness of such
    inferring for the truth.
    By means of such syllogisms (according to the side indicated in
    the section), the most diverse sorts of things can be proved, as it is
    said. The only thing required is to take up the medius terminus from
    which the transition to the desired determination can be made. Yet,
    with a different medius terminus, something else, even something
    opposite, may be proven. - The more concrete an object is, the more
    sides it has that inhere in it and can serve as medii termini. Which of
    these sides is more essential than the other must depend again on the
    sort of inferring that fixes upon the individual determinacy and can
    likewise easily find for it a side and a respect in terms of which it can
    be rendered important and necessarily valid.

Addition. As little as in the daily course of life one tends even to think of inference
at the level of understanding, it nonetheless plays its role incessantly there. Thus,
for example, in the civil dispute of law, it is the business of the advocates to make
the most favourable claim to a right for their parties. In a logical
respect, however, such a claim to a right is nothing other than a medius terminus.
The same also takes place in diplomatic negotiations if, for example, diverse powers
lay claim to one and the same land. In this connection, the right of inheritance, the
geographical location of the land, the descendancy and language of its inhabitants,
or any sort of other ground can be taken up as medius terminus. in the
second figure of the syllogism ((2) U -I -P}. This expresses the truth of
the first figure, [namely] that the mediation took place in the individuality
and accordingly is something contingent.

(I') This syllogism is equally contingent on account of the form of the
relation in it. According to the concept of the syllogism, the true is the
relation of differentiated entities, through a middle that
is their unity. The relations of the extremes to the middle (the so-called
premises, the major and the minor) are, however,
immediate relations.
This contradiction of the syllogism expresses itself again through an
infinite progression as the demand that each of the premises
likewise be proven by means of a syllogism; since this syllogism,
however, has two immediate premises of the same sort, this demand
then repeats itself and, indeed, as a demand constantly doubling
itself, ad infinitum.

What here (on account of the empirical importance) has been noted as
a deficiency of the syllogism, to which in this form absolute correctness
is ascribed, must of itself sublate itself in the
further determination of the syllogism. Here, within the sphere of the
concept as well as in the judgment, the opposite determinacy is not merely
in itself on hand, but instead it is posited, and hence, for the further
determination of the syllogism, it is only necessary to take up what is
posited each time by it itself.

By means of the immediate inference (I - P - U), the individual is
mediated with the universal and posited as universal in this conclusion.
By this means, the individual as subject, thus itself as universal,
is now the unity of the two extremes and the mediating factor, which results
The second figure joins the universal with the particular (i.e. the universal
that emerges from the previous conclusion is determined by the individ-
uality, and accordingly occupies the position of the immediate subject).
As a result, the universal is posited as particular, via the conclusion [of the
second figure], thus as the factor mediating the extremes, the positions of
which are now taken by the others in the third figure of the syllogism ((3)
P-U-I).
The so-called figures of the syllogism (Aristotle rightly acknowledges
only three of them; the fourth is a superficial, indeed fatuous,
addition of the moderns) are placed next to one another in the
standard treatment of them, without the slightest thought being
given to showing their necessity, even less their meaning and their
value. For this reason it is no wonder, if the figures have later been
treated as an empty formalism. They have, however, a very basic
[methodical] sense that rests upon the necessity that each moment as
a determination of the concept becomes itself the whole and the
mediating ground. - What determinations the sentences otherwise
have, whether they may be universal and so forth or negative,
in order to bring about a correct inference, this is a mechanical
investigation that has rightly come to be forgotten on account of
its concept-less mechanism and its lack of inner meaning. - One
can appeal least of all to Aristotle for the importance of such an
investigation and of syllogism at the level of understanding. To be
sure, he described these like countless other forms of the spirit and
nature, and he both investigated and presented their determinacy.
But in his metaphysical concepts as well as in the concepts of the
natural and the spiritual, he was far from intending to make the
form of the syllogism at the level of the understanding a foundation
and criterion, so far that probably not a single one of these concepts
would have been able to arise or be left standing if it were supposed
to be subjected to the laws of understanding. In the considerable
amount of descriptive and sensible detail that
Aristotle, after his fashion, brings together, the speculative concept is
invariably what dominates for him and he does not allow that
inferring at the level of mere understanding, the inferring that he
first outlined in so determinate a fashion, to enter into this sphere.

Addition. The objective sense of the figures of the syllogism is in general this,
that everything rational demonstrates itself in the form of a threefold syllogism
and, to be sure, in such a way that each of its members occupies equally the
position of an extreme and that of the mediating middle. This is expressly the
case with the three members of philosophical science, i.e. the logical idea, nature,
and spirit. First, nature here is the middle member that joins the others together.
Nature, this immediate totality, unfolds into the two extremes of the logical idea
and spirit. The spirit, however, is spirit only by being mediated by nature. Second,
then, the spirit, which we know as individual and active, is the middle,
and nature and the logical idea are the extremes. It is the spirit that recognizes
in nature the logical idea and elevates it to its essence. Third, the logical idea
is similarly the middle; it is the absolute substance of the spirit as of nature,
the universal, what pervades everything. These are the members of the absolute
syllogism.

Since each moment has run through the position of the middle and the
extremes, their determinate difference relative to one another has sublated
itself and, in this form where there is no difference between its moments, the
syllogism first has the external identity of the understanding, the equality,
as its relation the quantitative or mathematical syllogism. If
two things are equal to a third, then they are equal to one another.

Addition. The quantitative syllogism mentioned here surfaces familiarly in
mathematics as an axiom. It is customarily said of it, as of the other axioms,
that its content cannot be proven, but also that this proof is not needed since
it is immediately evident. Nevertheless, these mathematical axioms are in fact
nothing other than logical sentences that are to be derived, insofar as partic-
ular and determinate thoughts can be articulated in them, from universal and
self-determining thinking, a derivation which has to be considered then as their
proof. This is the case here with the quantitative syllogism, set up in mathematics
as an axiom that demonstrates itself to be the next result of the qualitative or
immediate syllogism. - The quantitative syllogism, moreover, is the completely
formless inference since in it the difference between the members, a difference
determined by the concept, is sublated. Which sentences here are supposed to be
premises depends upon external circumstances and, for this reason, in the appli-
cation of this syllogism one presupposes what already stands fast and is proven
elsewhere.

By this means, it has come about with respect to the form (1) that each
moment received the determination and position of the middle, hence
the whole in general, and with this it has lost the one-sidedness of its
abstraction in itself, (2) that the mediation (§ 185) has been
completed, also only in itself, namely, only as a circle of mediations that
mutually presuppose one another. In the first figure (I P U), the two
premises, 1- P and P - U, are still unmediated; the former being mediated
in the third, the latter in the second figure. But each of these two figures
equally presupposes the two other figures to mediate their premises.
In keeping with this, the mediating unity of the concept is no longer
to be posited only as an abstract particularity, but instead as the developed
unity of individuality and universality and, indeed, at first as the reflected
unity of these determinations, the individuality determined at the same time
as universality. This sort of middle yields the syllogism of reflection.

B. SYLLOGISM OF REFLECTION

The middle is in the first place (1) not alone the abstract, particular deter-
minacy of the subject, but instead at the same time as all individual concrete
subjects, to which that determinacy as only one among others accrues. As
such, the middle yields the syllogism of the set of all.
The major premise (the subject of which is the particular determinacy, the
terminus medius, as the set of all) presupposes the conclusion, of which it
is supposed to be the presupposition. It thus rests upon (2) induction, the
middle of which is the complete set of the individuals as such, a,
b, c, d, and so forth. Since, however, the immediate empirical individuality
is different from the universality and, for that reason, cannot ensure any
completeness, the induction rests upon (3) analogy, the middle of which
is an individual but in the sense of its essential universality, its genus or
essential determinacy. - The first syllogism refers, for its mediation, to
the second and the second to the third; but the latter equally demands a
universality or the individuality as genus after the forms of the external
relation of individuality and universality have been run through in the
figures of the syllogism of reflection.

By means of the syllogism of the set of all, some improvement is made
relative to the deficiency of the basic form of the inference at the level
of the understanding (pointed out in § 184). But the improvement is
only such that a new deficiency arises, namely, that the major premise
presupposes as an accordingly immediate sentence what was supposed to
be the conclusion. - 'All human beings are mortal, therefore Gajus is mortal',
'all metals are electric conductors, therefore, for example, copper is, too'. In
order to be able to assert those major premises that are supposed to express
the set of all of the immediate individuals and to be essentially empirical
sentences, it is required that already previously the sentences about the
individual Gajus, the individual copper are confirmed for themselves as
correct. - Everyone rightly notices not merely the pedantry, but the vapid
formalism of such syllogisms as 'all humans are mortal, but
now Gajus is human, and so forth'.

Addition. The syllogism with respect to the set of all refers to the syllogism
of induction in which the individuals form the middle that joins together the extremes.
If we say 'all metals are electric conductors', this is an empirical sentence
that results from testing undertaken with all individual metals. By this means, we
get the inference of induction, which has the following form: P-I U I I
Gold is metal, silver is metal; similarly, copper, lead, and so forth. This is the
major premise. Then comes the minor premise 'all these bodies are electric conduc-
tors', and from this results the conclusion that all metals are electric conductors.
Hence, here the individuality in the sense of the set of all is the binding factor.
This syllogism then likewise sends us on to another syllogism in turn. It has, as
its middle, the complete set of individuals. This presupposes that the observation
and experience be completed in a certain domain. But because it is a matter of
individualities here, this yields in turn the progression into infinity (I, I, I...).
In an induction the individuals can never be exhausted. If one says
'all metals', 'all plants', and so forth, this only means as much as 'all metals, all plants
with which one is familiar up to now'. Each induction is, therefore, imperfect.
One has, indeed, made this and that observation, one has made many observations,
but not all cases, not all individuals have been observed. It is this deficiency of
induction that leads to analogy. From the fact that things of a certain genus have a
certain property, it is inferred in the syllogism of analogy that other things of the
same genus have the same property. Thus, for example, it is a syllogism of analogy
if it is said: this law of motion has been found previously to hold for all planets;
hence, a newly discovered planet will probably move according to the same law.
In empirical sciences, analogy is rightly held in high regard and very important
results have been attained on this path. It is the instinct of reason that has the
presentiment that this or that empirically uncovered determination is grounded in
the inner nature or the genus of an object and whIch IS further based on this.
Incidentally, analogy can be more superficial or more rigorous. Suppose,
for example, it is said: 'Gaius, a human being, is a learned individual; Titus is also
a human being; hence, he is also likely to be learned: This is in any case a very
bad analogy and, indeed, for this reason, that for a human being to be learned is
not grounded without further ado in this, his genus. Nonetheless, .the same sort
of superficial analogies occur very frequently. Thus, for example, it is customarily
said: 'the Earth is a heavenly body and has inhabitants; the Moon is also a heavenly
body; hence, it is probably also inhabited'. This analogy is not a bit better than the
previously mentioned one. That the Earth has inhabitants does not rest merely on
the fact that it is a heavenly body; in addition, further conditions are also required,
such as, first and foremost, the fact that it is surrounded by an atmosphere, that
there is water on hand (something connected to that atmosphere), and so forth -
conditions that, as far as we know, are lacking in the case of the Moon. What is
called 'philosophy of nature' in the modern era consists to a great extent in a vapid
play with empty, external analogies, which are nonetheless supposed to count as
profound results. On account of this, philosophical consideration of nature has
fallen into a deserved disrepute.

C. SYLLOGISM OF NECESSITY

As far as the merely abstract determinations of this syllogism are concerned,
it has the universal as the middle term, just as the syllogism of reflection
has the individuality as the middle term,
the latter in terms of the second figure, the former in terms of the third;
the universal posited as essentially determined in itself.
(1) At first, the particular in the sense of the determinate genus or species is
the mediating determination - in the categorical syllogism;
(2) the individual in the sense of immediate being that is
both mediated and mediating in the hypothetical syllogism;
(3) the mediating universal is also posited as the totality of
its particularizations and as an individual particular,
 an exclusive individuality - in the disjunctive syllogism;
so that one and the same universal is in these determinations as
merely in forms of difference.

The syllogism has been taken in terms of the differences contained in it and
the universal result of the course of those differences is that
these differences and the concept's manner of being-outside-itself
sublate themselves in it.
Indeed, (1) each of the moments themselves has demonstrated itself to be
the totality of the moments and, hence, to be the entire syllogism;
thus, they are in themselves identical
(2) The negation of their differences and their mediation constitutes the
manner of being-for-itself, such that it is one and the same universal
that is in these forms and is accordingly also posited as their identity.
In this ideality of the moments, inferring acquires the determination
of essentially containing the negation of the determinacies by means of
which it runs its course and, with this, the determination of being a
mediation by way of sublating the mediation and a manner of joining
the subject together, not with another, but with the sublated other,
with itself.

Addition. In ordinary logic, the first part, forming the so-called 'doctrine of elements',
tends to come to a close with the treatment of the doctrine of the syllogism.
Following upon this is the so-called 'doctrine of method', as the second part,
in which it is supposed to be demonstrated how an entire body of scientific
knowledge is to be brought about, through application of the forms of thinking,
treated in the doctrine of elements, to the objects at hand. Logic at the level
of the understanding provides no further information about where these objects
come from and what sort of a connection it has in general with the thought of
objectivity. Thinking counts here as a merely subjective and formal activity, while
what is objective, in contrast to thinking, counts as something firm and on hand
for itself. This dualism, however, is not the true state of things and
it is a thoughtless procedure to take up the determinations of subjectivity and
objectivity without further ado and to refrain from inquiring into their origin.
Both, subjectivity as well as objectivity, are in any case thoughts and, to be sure,
determinate thoughts which have to demonstrate themselves in universal and self-
determining thinking. This has happened here first with respect to subjectivity.
We have come to recognize this, or the subjective concept, which contains in itself
the concept as such, the judgment, and the syllogism as the dialectical result of
the first two major stages of the logical idea, namely, the stages of being and of
essence. If it is said of the concept, 'it is subjective and only subjective', then this
is completely correct insofar as it is, to be sure, subjectivity itself But then, in
addition, no less subjective than the concept as such are the judgment and the
syllogism, determinations that in ordinary logic, next to the so-called 'laws of
thought' (the principles of identity, difference, and sufficient reason), form the
content of the so-called 'doctrine of elements'. Furthermore, this subjectivity with
the determinations of it mentioned here (the concept, the judgment, and the
syllogism) is not to be considered an empty framework which first has to acquire
its filling from without, through objects on hand for themselves. Instead it is the
subjectivity itself that, as dialectical, breaks through its limitation and by means
of the syllogism discloses itself to be objectivity.
